% from ega3-1-fr.tex
\oldpage[III]{128}

d'où on conclut que, pour tout sous-$A$-module $P$ de
$H^{n+1}(X,\mathfrak J^k\mathcal F)$, on a, dans le $S$-module
$\mathrm N$,
\[
\nu_m(S_mP)=\alpha_m(S_m)P.
\tag{4.1.7.19}\label{III.4.1.7.19-fr}
\]

\begin{lemma}[4.1.7.20]
\label{III.4.1.7.20-fr}
Il existe un entier $m>0$ tel que $\nu_m(\mathrm Q_{k+m})=0$ pour tout
$k\geq k_0$.
\end{lemma}

\begin{proof}
Prenons en effet pour $m$ un multiple de $h$ qui soit $\geq r$~; comme
$\mathrm Q_{k+m}=S_m\mathrm Q_k$ pour $k\geq k_0$, on a en vertu de
(\hyperref[III.4.1.7.19-fr]{4.1.7.19}) et
(\hyperref[III.4.1.7.15-fr]{4.1.7.15})
\[
\nu_m(\mathrm Q_{k+m})=\alpha_m(S_m)\mathrm Q_k
\subset\alpha_r(S_r)\mathrm Q_k=0.
\]
\end{proof}

\begin{env}[4.1.7.21]
\label{III.4.1.7.21-fr}
Remarquons que du diagramme commutatif
\[
\xymatrix@C=2.5em{
H^n(X,\mathcal F)\ar[r]&H^n(X,\mathcal F_k)\ar[r]^-\partial&
H^{n+1}(X,\mathfrak J^{k+1}\mathcal F)\ar[r]&H^{n+1}(X,\mathcal F)\\
H^n(X,\mathcal F)\ar[r]\ar[u]&H^n(X,\mathcal F_{k+m})\ar[r]^-\partial\ar[u]&
H^{n+1}(X,\mathfrak J^{k+m+1}\mathcal F)\ar[r]\ar[u]&
H^{n+1}(X,\mathcal F)\ar[u]
}
\]
provenant lui-même du diagramme commutatif
\[
\xymatrix{
0\ar[r]&\mathfrak J^{k+1}\mathcal F\ar[r]&\mathcal F\ar[r]&
\mathcal F_k\ar[r]&0\\
0\ar[r]&\mathfrak J^{k+m+1}\mathcal F\ar[r]\ar[u]&
\mathcal F\ar[r]\ar[u]&\mathcal F_{k+m}\ar[r]\ar[u]&0
}
\]
où les flèches verticales sont les applications canoniques, on déduit un
diagramme commutatif
\[
\xymatrix@C=2.2em{
0\ar[r]&\mathrm R_k\ar[r]&\mathrm H\ar[r]&\mathrm H_k\ar[r]&
\mathrm Q_k\ar[r]&0\\
0\ar[r]&\mathrm R_{k+m}\ar[r]\ar[u]&\mathrm H\ar[r]\ar@{=}[u]&
\mathrm H_{k+m}\ar[r]\ar[u]&\mathrm Q_{k+m}\ar[r]\ar[u]^{\nu_m}&0
}
\]
où les lignes sont exactes. Comme la dernière flèche verticale est nulle
pour $k\geq k_0$ (\hyperref[III.4.1.7.20-fr]{4.1.7.20}), l'image de
$\mathrm H_{k+m}$ dans $\mathrm H_k$ est contenue dans
$\operatorname{Ker}(\mathrm H_k\to\mathrm Q_k)
=\operatorname{Im}(\mathrm H\to\mathrm H_k)$, mais par ailleurs elle
contient $\operatorname{Im}(\mathrm H\to\mathrm H_k)$ par la commutativité
du diagramme, donc elle lui est égale~; il en est donc de même des images
dans $\mathrm H_k$ des $\mathrm H_{k'}$ pour $k'\geq k+m$, ce qui démontre
la condition (ML) pour le système projectif $(\mathrm H_k)_{k\geq0}$. En
outre, pour tout ouvert affine $U$ de $X$, on a
$H^i(U,\mathcal F_k)=0$ pour $i>0$
(\hyperref[III.1.3.1-fr]{1.3.1}), et pour $m>0$, l'application
$H^0(U,\mathcal F_{k+m})\to H^0(U,\mathcal F_k)$ est surjective
(I, 1.3.9). On peut donc appli-

